\Chapter{74}

\Verse{1}
{
    \zA{话说平儿听迎春说了，正自好笑，忽见宝玉也来了。}
}
{
    \eA{[English source not present in the checked-in Bencraft/Joly files.]}
}
{
}

\Verse{2}
{
    \zA{原来管厨房柳家媳妇的妹 子也因放头开赌，得了不是，因这园中有素和柳家的不好的，便又告出柳家的来， 说和他妹子是伙计，赚了平分。}
}
{
    \eA{[English source not present in the checked-in Bencraft/Joly files.]}
}
{
}

\Verse{3}
{
    \zA{因此凤姐要治柳家之罪。}
}
{
    \eA{[English source not present in the checked-in Bencraft/Joly files.]}
}
{
}

\Verse{4}
{
    \zA{那柳家的听得此言，便慌 了手脚，因思素与怡红院的人最为深厚，故走来悄悄的央求晴雯芳官等人，转告诉 了宝玉。}
}
{
    \eA{[English source not present in the checked-in Bencraft/Joly files.]}
}
{
}

\Verse{5}
{
    \zA{宝玉因思内中迎春的嬷嬷也现有此罪，不若来约同迎春去讨情，比自己独 去单为柳家的说情又更妥当，故此前来。}
}
{
    \eA{[English source not present in the checked-in Bencraft/Joly files.]}
}
{
}

\Verse{6}
{
    \zA{忽见许多人在此，见他来时，都问道：“你 的病可好了？}
}
{
    \eA{[English source not present in the checked-in Bencraft/Joly files.]}
}
{
}

\Verse{7}
{
    \zA{跑来做什么？”}
}
{
    \eA{[English source not present in the checked-in Bencraft/Joly files.]}
}
{
}

\Verse{8}
{
    \zA{宝玉不便说出讨情一事，只说：“来看二姐姐。”}
}
{
    \eA{[English source not present in the checked-in Bencraft/Joly files.]}
}
{
}

\Verse{9}
{
    \zA{当 下众人也不在意，且说些闲话。}
}
{
    \eA{[English source not present in the checked-in Bencraft/Joly files.]}
}
{
}

\Verse{10}
{
    \zA{平儿便出去办累金凤一事。}
}
{
    \eA{[English source not present in the checked-in Bencraft/Joly files.]}
}
{
}

\Verse{11}
{
    \zA{那玉柱儿媳妇紧跟在后，口内百般央求，只说：“姑 娘好歹口内超生，我横竖去赎了来。”}
}
{
    \eA{[English source not present in the checked-in Bencraft/Joly files.]}
}
{
}

\Verse{12}
{
    \zA{平儿笑道：“你迟也赎，早也赎，‘既有今 日，何必当初’。}
}
{
    \eA{[English source not present in the checked-in Bencraft/Joly files.]}
}
{
}

\Verse{13}
{
    \zA{你的意思得过就过，既这么样，我也不好意思告诉人。}
}
{
    \eA{[English source not present in the checked-in Bencraft/Joly files.]}
}
{
}

\Verse{14}
{
    \zA{趁早儿取 了来，交给我，一字不提。”}
}
{
    \eA{[English source not present in the checked-in Bencraft/Joly files.]}
}
{
}

\Verse{15}
{
    \zA{玉柱儿媳妇听说，方放下心来，就拜谢，又说：“姑 娘自去贵干。}
}
{
    \eA{[English source not present in the checked-in Bencraft/Joly files.]}
}
{
}

\Verse{16}
{
    \zA{赶晚赎了来，先回了姑娘再送去如何？”}
}
{
    \eA{[English source not present in the checked-in Bencraft/Joly files.]}
}
{
}

\Verse{17}
{
    \zA{平儿道：“赶晚不来，可别 怨我！”}
}
{
    \eA{[English source not present in the checked-in Bencraft/Joly files.]}
}
{
}

\Verse{18}
{
    \zA{说毕，二人方分路各自散了。}
}
{
    \eA{[English source not present in the checked-in Bencraft/Joly files.]}
}
{
}

\Verse{19}
{
    \zA{平儿到房，凤姐问他：“三姑娘叫你做什么？”}
}
{
    \eA{[English source not present in the checked-in Bencraft/Joly files.]}
}
{
}

\Verse{20}
{
    \zA{平儿笑道：“三姑娘怕奶奶生气，叫我劝着奶奶些，问奶奶这两天可吃些什么？”}
}
{
    \eA{[English source not present in the checked-in Bencraft/Joly files.]}
}
{
}

\Verse{21}
{
    \zA{凤姐笑道：“倒是他还惦记我。}
}
{
    \eA{[English source not present in the checked-in Bencraft/Joly files.]}
}
{
}

\Verse{22}
{
    \zA{刚才又出来了一件事：有人来告柳二媳妇和他妹子 通同开局，凡妹子所为都是他作主。}
}
{
    \eA{[English source not present in the checked-in Bencraft/Joly files.]}
}
{
}

\Verse{23}
{
    \zA{我想你素日肯劝我多一事不如少一事，自己保 养保养也是好的。}
}
{
    \eA{[English source not present in the checked-in Bencraft/Joly files.]}
}
{
}

\Verse{24}
{
    \zA{我因听不进去，果然应了，先把太太得罪了，而且反赚了一场病。}
}
{
    \eA{[English source not present in the checked-in Bencraft/Joly files.]}
}
{
}

\Verse{25}
{
    \zA{如今我也看破了，随他们闹去罢，横竖还有许多人呢。}
}
{
    \eA{[English source not present in the checked-in Bencraft/Joly files.]}
}
{
}

\Verse{26}
{
    \zA{我白操一会子心，倒惹的万 人咒骂，不如且自家养养病。}
}
{
    \eA{[English source not present in the checked-in Bencraft/Joly files.]}
}
{
}

\Verse{27}
{
    \zA{就是病好了，我也会做好好先生，得乐且乐，得笑且 笑，一概是非都凭他们去罢，所以我只答应着‘知道了’。”}
}
{
    \eA{[English source not present in the checked-in Bencraft/Joly files.]}
}
{
}

\Verse{28}
{
    \zA{平儿笑道：“奶奶果 然如此，那就是我们的造化了。”}
}
{
    \eA{[English source not present in the checked-in Bencraft/Joly files.]}
}
{
}

\Verse{29}
{
    \zA{一语未了，只见贾琏进来，拍手叹气道：“好好的又生事!前儿我和鸳鸯借当， 那边太太怎么知道了？}
}
{
    \eA{[English source not present in the checked-in Bencraft/Joly files.]}
}
{
}

\Verse{30}
{
    \zA{刚才太太叫过我去，叫我不管那里先借二百银子，做八月十 五节下使用。}
}
{
    \eA{[English source not present in the checked-in Bencraft/Joly files.]}
}
{
}

\Verse{31}
{
    \zA{我回没处借，太太就说：‘你没有钱就有地方挪移，我白和你商量， 你就搪塞我!你就没地方儿!前儿一千银子的当是那里的？}
}
{
    \eA{[English source not present in the checked-in Bencraft/Joly files.]}
}
{
}

\Verse{32}
{
    \zA{连老太太的东西你都有神 通弄出来，这会二百银子你就这样难。}
}
{
    \eA{[English source not present in the checked-in Bencraft/Joly files.]}
}
{
}

\Verse{33}
{
    \zA{亏我没和别人说去！’}
}
{
    \eA{[English source not present in the checked-in Bencraft/Joly files.]}
}
{
}

\Verse{34}
{
    \zA{我想太太分明不短， 何苦来又寻事奈何人！”}
}
{
    \eA{[English source not present in the checked-in Bencraft/Joly files.]}
}
{
}

\Verse{35}
{
    \zA{凤姐儿道：“那日并没个外人，谁走了这个消息？”}
}
{
    \eA{[English source not present in the checked-in Bencraft/Joly files.]}
}
{
}

\Verse{36}
{
    \zA{平儿 听了，也细想那日有谁在此，想了半日，笑道：“是了。}
}
{
    \eA{[English source not present in the checked-in Bencraft/Joly files.]}
}
{
}

\Verse{37}
{
    \zA{那日说话时没人，就只晚 上送东西来的时候儿，老太太那边傻大姐的娘可巧来送浆洗衣裳，他在下房里坐了 一会子，看见一大箱子东西，自然要问。}
}
{
    \eA{[English source not present in the checked-in Bencraft/Joly files.]}
}
{
}

\Verse{38}
{
    \zA{必是丫头们不知道，说出来了，也未可知。”}
}
{
    \eA{[English source not present in the checked-in Bencraft/Joly files.]}
}
{
}

\Verse{39}
{
    \zA{因此便唤了几个小丫头来问：“那日谁告诉傻大姐的娘了？”}
}
{
    \eA{[English source not present in the checked-in Bencraft/Joly files.]}
}
{
}

\Verse{40}
{
    \zA{众小丫头慌了，都跪 下赌神发誓说：“自来也没敢多说一句话。}
}
{
    \eA{[English source not present in the checked-in Bencraft/Joly files.]}
}
{
}

\Verse{41}
{
    \zA{有人凡问什么，都答应不知道，这事如 何敢说！”}
}
{
    \eA{[English source not present in the checked-in Bencraft/Joly files.]}
}
{
}

\Verse{42}
{
    \zA{凤姐详情度理，说：“他们必不敢多说一句话，倒别委屈了他们。}
}
{
    \eA{[English source not present in the checked-in Bencraft/Joly files.]}
}
{
}

\Verse{43}
{
    \zA{如今 把这事靠后，且把太太打发了去要紧。}
}
{
    \eA{[English source not present in the checked-in Bencraft/Joly files.]}
}
{
}

\Verse{44}
{
    \zA{宁可咱们短些，别又讨没意思。”}
}
{
    \eA{[English source not present in the checked-in Bencraft/Joly files.]}
}
{
}

\Verse{45}
{
    \zA{因叫平儿： “把我的金首饰再去押二百银子来，送去完事。”}
}
{
    \eA{[English source not present in the checked-in Bencraft/Joly files.]}
}
{
}

\Verse{46}
{
    \zA{贾琏道：“索性多押二百，咱们 也要使呢。”}
}
{
    \eA{[English source not present in the checked-in Bencraft/Joly files.]}
}
{
}

\Verse{47}
{
    \zA{凤姐道：“很不必，我没处使。}
}
{
    \eA{[English source not present in the checked-in Bencraft/Joly files.]}
}
{
}

\Verse{48}
{
    \zA{这不知还指那一项赎呢。”}
}
{
    \eA{[English source not present in the checked-in Bencraft/Joly files.]}
}
{
}

\Verse{49}
{
    \zA{平儿拿了 去，吩咐旺儿媳妇领去。}
}
{
    \eA{[English source not present in the checked-in Bencraft/Joly files.]}
}
{
}

\Verse{50}
{
    \zA{不一时拿了银子来，贾琏亲自送去，不在话下。}
}
{
    \eA{[English source not present in the checked-in Bencraft/Joly files.]}
}
{
}

\Verse{51}
{
    \zA{这里凤姐和平儿猜疑走风的人：“反叫鸳鸯受累，岂不是咱们之过！”}
}
{
    \eA{[English source not present in the checked-in Bencraft/Joly files.]}
}
{
}

\Verse{52}
{
    \zA{正在胡 想，人报：“太太来了。”}
}
{
    \eA{[English source not present in the checked-in Bencraft/Joly files.]}
}
{
}

\Verse{53}
{
    \zA{凤姐听了诧异，不知何事，遂与平儿等忙迎出来。}
}
{
    \eA{[English source not present in the checked-in Bencraft/Joly files.]}
}
{
}

\Verse{54}
{
    \zA{只见 王夫人气色更变，只带一个贴己小丫头走来，一语不发，走至里间坐下。}
}
{
    \eA{[English source not present in the checked-in Bencraft/Joly files.]}
}
{
}

\Verse{55}
{
    \zA{凤姐忙捧 茶，因陪笑问道：“太太今日高兴，到这里逛逛？”}
}
{
    \eA{[English source not present in the checked-in Bencraft/Joly files.]}
}
{
}

\Verse{56}
{
    \zA{王夫人喝命：“平儿出去！”}
}
{
    \eA{[English source not present in the checked-in Bencraft/Joly files.]}
}
{
}

\Verse{57}
{
    \zA{平儿见了这般，不知怎么了，忙应了一声，带着众小丫头一齐出去，在房门外站住。}
}
{
    \eA{[English source not present in the checked-in Bencraft/Joly files.]}
}
{
}

\Verse{58}
{
    \zA{一面将房门掩了，自己坐在台阶上，所有的人一个不许进去。}
}
{
    \eA{[English source not present in the checked-in Bencraft/Joly files.]}
}
{
}

\Verse{59}
{
    \zA{凤姐也着了慌，不知 有何事。}
}
{
    \eA{[English source not present in the checked-in Bencraft/Joly files.]}
}
{
}

\Verse{60}
{
    \zA{只见王夫人含着泪，从袖里扔出一个香袋来，说：“你瞧！”}
}
{
    \eA{[English source not present in the checked-in Bencraft/Joly files.]}
}
{
}

\Verse{61}
{
    \zA{凤姐忙拾起 一看，见是十锦春意香袋，也吓了一跳，忙问：“太太从那里得来？”}
}
{
    \eA{[English source not present in the checked-in Bencraft/Joly files.]}
}
{
}

\Verse{62}
{
    \zA{王夫人见问， 越发泪如雨下，颤声说道：“我从那里得来？}
}
{
    \eA{[English source not present in the checked-in Bencraft/Joly files.]}
}
{
}

\Verse{63}
{
    \zA{我天天坐在井里!想你是个细心人，所 以我才偷空儿，谁知你也和我一样!这样东西，大天白日，明摆在园里山石上，被 老太太的丫头拾着。}
}
{
    \eA{[English source not present in the checked-in Bencraft/Joly files.]}
}
{
}

\Verse{64}
{
    \zA{不亏你婆婆看见，早已送到老太太跟前去了。}
}
{
    \eA{[English source not present in the checked-in Bencraft/Joly files.]}
}
{
}

\Verse{65}
{
    \zA{我且问你：这个 东西如何丢在那里？”}
}
{
    \eA{[English source not present in the checked-in Bencraft/Joly files.]}
}
{
}

\Verse{66}
{
    \zA{凤姐听得，也更了颜色，忙问：“太太怎么知道是我的？”}
}
{
    \eA{[English source not present in the checked-in Bencraft/Joly files.]}
}
{
}

\Verse{67}
{
    \zA{王夫人又哭又叹道：“你反问我？}
}
{
    \eA{[English source not present in the checked-in Bencraft/Joly files.]}
}
{
}

\Verse{68}
{
    \zA{你想，一家子除了你们小夫小妻，馀者老婆子们， 要这个何用？}
}
{
    \eA{[English source not present in the checked-in Bencraft/Joly files.]}
}
{
}

\Verse{69}
{
    \zA{女孩子们是从那里得来？}
}
{
    \eA{[English source not present in the checked-in Bencraft/Joly files.]}
}
{
}

\Verse{70}
{
    \zA{自然是那琏儿不长进下流种子那里弄来的。}
}
{
    \eA{[English source not present in the checked-in Bencraft/Joly files.]}
}
{
}

\Verse{71}
{
    \zA{你 们又和气，当作一件玩意儿。}
}
{
    \eA{[English source not present in the checked-in Bencraft/Joly files.]}
}
{
}

\Verse{72}
{
    \zA{年轻的人，儿女闺房私意是有的，你还和我赖!幸而 园内上下人还不解事，尚未拣得，倘或丫头们拣着，你姊妹看见，这还了得？}
}
{
    \eA{[English source not present in the checked-in Bencraft/Joly files.]}
}
{
}

\Verse{73}
{
    \zA{不然， 有那小丫头们拣着出去，说是园内拣的，外人知道，这性命脸面要也不要？”}
}
{
    \eA{[English source not present in the checked-in Bencraft/Joly files.]}
}
{
}

\Verse{74}
{
    \zA{凤姐听说，又急又愧，登时紫胀了面皮，便挨着炕沿双膝跪下，也含泪诉道： “太太说的固然有理，我也不敢辩。}
}
{
    \eA{[English source not present in the checked-in Bencraft/Joly files.]}
}
{
}

\Verse{75}
{
    \zA{但我并无这样东西，其中还要求太太细想：这 香袋儿是外头仿着内工绣的，连穗子一概都是市卖的东西。}
}
{
    \eA{[English source not present in the checked-in Bencraft/Joly files.]}
}
{
}

\Verse{76}
{
    \zA{我虽年轻不尊重，也不 肯要这样东西。}
}
{
    \eA{[English source not present in the checked-in Bencraft/Joly files.]}
}
{
}

\Verse{77}
{
    \zA{再者，这也不是常带着的，我纵然有，也只好在私处搁着，焉肯在 身上常带，各处逛去？}
}
{
    \eA{[English source not present in the checked-in Bencraft/Joly files.]}
}
{
}

\Verse{78}
{
    \zA{况且又在园里去，个个姊妹，我们都肯拉拉扯扯，倘或露出 来，不但在姊妹前看见，就是奴才看见，我有什么意思？}
}
{
    \eA{[English source not present in the checked-in Bencraft/Joly files.]}
}
{
}

\Verse{79}
{
    \zA{三则论主子内我是年轻媳 妇，算起来，奴才比我更年轻的又不止一个了，况且他们也常在园走动，焉知不是 他们掉的？}
}
{
    \eA{[English source not present in the checked-in Bencraft/Joly files.]}
}
{
}

\Verse{80}
{
    \zA{再者，除我常在园里，还有那边太太常带过几个小姨娘来，嫣红翠云那 几个人也都是年轻的人，他们更该有这个了。}
}
{
    \eA{[English source not present in the checked-in Bencraft/Joly files.]}
}
{
}

\Verse{81}
{
    \zA{还有那边珍大嫂子，他也不算很老， 也常带过佩凤他们来，又焉知又不是他们的？}
}
{
    \eA{[English source not present in the checked-in Bencraft/Joly files.]}
}
{
}

\Verse{82}
{
    \zA{况且园内丫头也多，保不住都是正经 的。}
}
{
    \eA{[English source not present in the checked-in Bencraft/Joly files.]}
}
{
}

\Verse{83}
{
    \zA{或者年纪大些的知道了人事，一刻查问不到，偷出去了，或借着因由合二门上 小么儿们打牙撂嘴儿，外头得了来的，也未可知。}
}
{
    \eA{[English source not present in the checked-in Bencraft/Joly files.]}
}
{
}

\Verse{84}
{
    \zA{不但我没此事，就连平儿，我也 可以下保的：太太请细想。”}
}
{
    \eA{[English source not present in the checked-in Bencraft/Joly files.]}
}
{
}

\Verse{85}
{
    \zA{王夫人听了这一席话，很近情理，因叹道：“你起来。}
}
{
    \eA{[English source not present in the checked-in Bencraft/Joly files.]}
}
{
}

\Verse{86}
{
    \zA{我也知道你是大家子的 姑娘出身，不至这样轻薄，不过我气激你的话。}
}
{
    \eA{[English source not present in the checked-in Bencraft/Joly files.]}
}
{
}

\Verse{87}
{
    \zA{但只如今且怎么处？}
}
{
    \eA{[English source not present in the checked-in Bencraft/Joly files.]}
}
{
}

\Verse{88}
{
    \zA{你婆婆才打发 人封了这个给我瞧，把我气了个死。”}
}
{
    \eA{[English source not present in the checked-in Bencraft/Joly files.]}
}
{
}

\Verse{89}
{
    \zA{凤姐道：“太太快别生气。}
}
{
    \eA{[English source not present in the checked-in Bencraft/Joly files.]}
}
{
}

\Verse{90}
{
    \zA{若被众人觉察了， 保不定老太太不知道。}
}
{
    \eA{[English source not present in the checked-in Bencraft/Joly files.]}
}
{
}

\Verse{91}
{
    \zA{且平心静气，暗暗访察，才能得这个实在；纵然访不着，外 人也不能知道。}
}
{
    \eA{[English source not present in the checked-in Bencraft/Joly files.]}
}
{
}

\Verse{92}
{
    \zA{如今惟有趁着赌钱的因由革了许多人这空儿，把周瑞媳妇、旺儿媳 妇等四五个贴近不能走话的人，安插在园里，以查赌为由。}
}
{
    \eA{[English source not present in the checked-in Bencraft/Joly files.]}
}
{
}

\Verse{93}
{
    \zA{再如今他们的丫头也太 多了，保不住人大心大，生事作耗，等闹出来，反悔之不及。}
}
{
    \eA{[English source not present in the checked-in Bencraft/Joly files.]}
}
{
}

\Verse{94}
{
    \zA{如今若无故裁革，不 但姑娘们委屈，就连太太和我也过不去。}
}
{
    \eA{[English source not present in the checked-in Bencraft/Joly files.]}
}
{
}

\Verse{95}
{
    \zA{不如趁着这个机会，以后凡年纪大些的， 或有些磨牙难缠的，拿个错儿撵出去，配了人：一则保的住没有别事，二则也可省 些用度。}
}
{
    \eA{[English source not present in the checked-in Bencraft/Joly files.]}
}
{
}

\Verse{96}
{
    \zA{太太想我这话如何？”}
}
{
    \eA{[English source not present in the checked-in Bencraft/Joly files.]}
}
{
}

\Verse{97}
{
    \zA{王夫人叹道：“你说的何尝不是。}
}
{
    \eA{[English source not present in the checked-in Bencraft/Joly files.]}
}
{
}

\Verse{98}
{
    \zA{但从公细想，你 这几个姊妹，每人只有两三个丫头像人，馀者竟是小鬼儿似的。}
}
{
    \eA{[English source not present in the checked-in Bencraft/Joly files.]}
}
{
}

\Verse{99}
{
    \zA{如今再去了，不但 我心里不忍，只怕老太太未必就依。}
}
{
    \eA{[English source not present in the checked-in Bencraft/Joly files.]}
}
{
}

\Verse{100}
{
    \zA{虽然艰难，也还穷不至此。}
}
{
    \eA{[English source not present in the checked-in Bencraft/Joly files.]}
}
{
}

\Verse{101}
{
    \zA{我虽没受过大荣华， 比你们是强些，如今宁可省我些，别委屈了他们。}
}
{
    \eA{[English source not present in the checked-in Bencraft/Joly files.]}
}
{
}

\Verse{102}
{
    \zA{你如今且叫人传周瑞家的等人进 来，就吩咐他们快快暗访这事要紧！”}
}
{
    \eA{[English source not present in the checked-in Bencraft/Joly files.]}
}
{
}

\Verse{103}
{
    \zA{凤姐即唤平儿进来，吩咐出去。}
}
{
    \eA{[English source not present in the checked-in Bencraft/Joly files.]}
}
{
}

\Verse{104}
{
    \zA{一时，周瑞家的与吴兴家的、郑华家的、来旺 家的、来喜家的现在五家陪房进来。}
}
{
    \eA{[English source not present in the checked-in Bencraft/Joly files.]}
}
{
}

\Verse{105}
{
    \zA{王夫人正嫌人少，不能勘察，忽见邢夫人的陪 房王善保家的走来，正是方才是他送香袋来的。}
}
{
    \eA{[English source not present in the checked-in Bencraft/Joly files.]}
}
{
}

\Verse{106}
{
    \zA{王夫人向来看视邢夫人之得力心腹 人等原无二意，今见他来打听此事，便向他说：“你去回了太太，也进园来照管照 管，比别人强些。”}
}
{
    \eA{[English source not present in the checked-in Bencraft/Joly files.]}
}
{
}

\Verse{107}
{
    \zA{王善保家的因素日进园去，那些丫鬟们不大趋奉他，他心里不 自在，要寻他们的故事又寻不着，恰好生出这件事来，以为得了把柄；又听王夫人
        委托他，正碰在心坎上，道：“这个容易。}
}
{
    \eA{[English source not present in the checked-in Bencraft/Joly files.]}
}
{
}

\Verse{108}
{
    \zA{不是奴才多话，论理这事该早严紧些的。}
}
{
    \eA{[English source not present in the checked-in Bencraft/Joly files.]}
}
{
}

\Verse{109}
{
    \zA{太太也不大往园里去，这些女孩子们，一个个倒像受了诰封似的，他们就成了千金 小姐了。}
}
{
    \eA{[English source not present in the checked-in Bencraft/Joly files.]}
}
{
}

\Verse{110}
{
    \zA{闹下天来，谁敢哼一声儿。}
}
{
    \eA{[English source not present in the checked-in Bencraft/Joly files.]}
}
{
}

\Verse{111}
{
    \zA{不然，就调唆姑娘们，说欺负了姑娘们了，谁 还耽得起！”}
}
{
    \eA{[English source not present in the checked-in Bencraft/Joly files.]}
}
{
}

\Verse{112}
{
    \zA{王夫人点头道：“跟姑娘们的丫头比别的娇贵些，这也是常情。”}
}
{
    \eA{[English source not present in the checked-in Bencraft/Joly files.]}
}
{
}

\Verse{113}
{
    \zA{王 善保家的道：“别的还罢了，太太不知，头一个是宝玉屋里的晴雯那丫头，仗着他 的模样儿比别人标致些，又长了一张巧嘴，天天打扮的像个西施样子，在人跟前能
        说惯道，抓尖要强。}
}
{
    \eA{[English source not present in the checked-in Bencraft/Joly files.]}
}
{
}

\Verse{114}
{
    \zA{一句话不投机，他就立起两只眼睛来骂人。}
}
{
    \eA{[English source not present in the checked-in Bencraft/Joly files.]}
}
{
}

\Verse{115}
{
    \zA{妖妖调调，大不成 个体统。”}
}
{
    \eA{[English source not present in the checked-in Bencraft/Joly files.]}
}
{
}

\Verse{116}
{
    \zA{王夫人听了这话，猛然触动往事，便问凤姐道：“上次我们跟了老太太 进园逛去，有一个水蛇腰，削肩膀儿，眉眼又有些像你林妹妹的，正在那里骂小丫
        头，我心里很看不上那狂样子。}
}
{
    \eA{[English source not present in the checked-in Bencraft/Joly files.]}
}
{
}

\Verse{117}
{
    \zA{因同老太太走，我不曾说他；后来要问是谁，偏又 忘了。}
}
{
    \eA{[English source not present in the checked-in Bencraft/Joly files.]}
}
{
}

\Verse{118}
{
    \zA{今日对了槛儿，这丫头想必就是他了？”}
}
{
    \eA{[English source not present in the checked-in Bencraft/Joly files.]}
}
{
}

\Verse{119}
{
    \zA{凤姐道：“若论这些丫头们，共总 比起来，都没晴雯长得好。}
}
{
    \eA{[English source not present in the checked-in Bencraft/Joly files.]}
}
{
}

\Verse{120}
{
    \zA{论举止言语，他原轻薄些。}
}
{
    \eA{[English source not present in the checked-in Bencraft/Joly files.]}
}
{
}

\Verse{121}
{
    \zA{方才太太说的倒很像他，我 也忘了那日的事，不敢混说。”}
}
{
    \eA{[English source not present in the checked-in Bencraft/Joly files.]}
}
{
}

\Verse{122}
{
    \zA{王善保家的便道：“不用这样，此刻不难叫了他来， 太太瞧瞧。”}
}
{
    \eA{[English source not present in the checked-in Bencraft/Joly files.]}
}
{
}

\Verse{123}
{
    \zA{王夫人道：“宝玉屋里常见我的，只有袭人麝月，这两个笨笨的倒好。}
}
{
    \eA{[English source not present in the checked-in Bencraft/Joly files.]}
}
{
}

\Verse{124}
{
    \zA{要有这个，他自然不敢来见我呀。}
}
{
    \eA{[English source not present in the checked-in Bencraft/Joly files.]}
}
{
}

\Verse{125}
{
    \zA{我一生最嫌这样的人，且又出来这个事。}
}
{
    \eA{[English source not present in the checked-in Bencraft/Joly files.]}
}
{
}

\Verse{126}
{
    \zA{好好的 宝玉倘或叫这蹄子勾引坏了，那还了得。”}
}
{
    \eA{[English source not present in the checked-in Bencraft/Joly files.]}
}
{
}

\Verse{127}
{
    \zA{因叫自己的丫头来，吩咐他道：“你去， 只说我有话问他，留下袭人麝月伏侍宝玉，不必来；有一个晴雯最伶俐，叫他即刻 快来。}
}
{
    \eA{[English source not present in the checked-in Bencraft/Joly files.]}
}
{
}

\Verse{128}
{
    \zA{你不许和他说什么！”}
}
{
    \eA{[English source not present in the checked-in Bencraft/Joly files.]}
}
{
}

\Verse{129}
{
    \zA{小丫头答应了，走入怡红院，正值晴雯身上不好，睡中觉才起来，发闷呢，听 如此说，只得跟了他来。}
}
{
    \eA{[English source not present in the checked-in Bencraft/Joly files.]}
}
{
}

\Verse{130}
{
    \zA{素日晴雯不敢出头，因连日不自在，并没十分妆饰，自为 无碍。}
}
{
    \eA{[English source not present in the checked-in Bencraft/Joly files.]}
}
{
}

\Verse{131}
{
    \zA{及到了凤姐房中，王夫人一见他钗？}
}
{
    \eA{[English source not present in the checked-in Bencraft/Joly files.]}
}
{
}

\Verse{132}
{
    \zA{鬓松，衫垂带褪，大有春睡捧心之态， 而且形容面貌恰是上月的那人，不觉勾起方才的火来。}
}
{
    \eA{[English source not present in the checked-in Bencraft/Joly files.]}
}
{
}

\Verse{133}
{
    \zA{王夫人便冷笑道：“好个美 人儿，真像个‘病西施’了。}
}
{
    \eA{[English source not present in the checked-in Bencraft/Joly files.]}
}
{
}

\Verse{134}
{
    \zA{你天天作这轻狂样儿给谁看!你干的事，打量我不知 道呢。}
}
{
    \eA{[English source not present in the checked-in Bencraft/Joly files.]}
}
{
}

\Verse{135}
{
    \zA{我且放着你，自然明儿揭你的皮!――宝玉今日可好些？”}
}
{
    \eA{[English source not present in the checked-in Bencraft/Joly files.]}
}
{
}

\Verse{136}
{
    \zA{晴雯一听如此说， 心内大异，便知有人暗算了他，虽然着恼，只不敢作声。}
}
{
    \eA{[English source not present in the checked-in Bencraft/Joly files.]}
}
{
}

\Verse{137}
{
    \zA{他本是个聪明过顶的人， 见问宝玉可好些，他便不肯以实话答应，忙跪下回道：“我不大到宝玉房里去，又
        不常和宝玉在一处，好歹我不能知，那都是袭人合麝月两个人的事，太太问他们。”}
}
{
    \eA{[English source not present in the checked-in Bencraft/Joly files.]}
}
{
}

\Verse{138}
{
    \zA{王夫人道：“这就该打嘴。}
}
{
    \eA{[English source not present in the checked-in Bencraft/Joly files.]}
}
{
}

\Verse{139}
{
    \zA{你难道是死人？}
}
{
    \eA{[English source not present in the checked-in Bencraft/Joly files.]}
}
{
}

\Verse{140}
{
    \zA{要你们做什么？”}
}
{
    \eA{[English source not present in the checked-in Bencraft/Joly files.]}
}
{
}

\Verse{141}
{
    \zA{晴雯道：“我原是跟 老太太的人，因老太太说园里空大，人少，宝玉害怕，所以拨了我去外间屋里上夜， 不过看屋子。}
}
{
    \eA{[English source not present in the checked-in Bencraft/Joly files.]}
}
{
}

\Verse{142}
{
    \zA{我原回过我笨，不能伏侍，老太太骂了我，‘又不叫你管他的事，要 伶俐的做什么？’}
}
{
    \eA{[English source not present in the checked-in Bencraft/Joly files.]}
}
{
}

\Verse{143}
{
    \zA{我听了不敢不去，才去的。}
}
{
    \eA{[English source not present in the checked-in Bencraft/Joly files.]}
}
{
}

\Verse{144}
{
    \zA{不过十天半月之内，宝玉叫着了，答 应几句话，就散了。}
}
{
    \eA{[English source not present in the checked-in Bencraft/Joly files.]}
}
{
}

\Verse{145}
{
    \zA{至于宝玉的饮食起居，上一层有老奶奶老妈妈们，下一层有袭 人、麝月、秋纹几个人。}
}
{
    \eA{[English source not present in the checked-in Bencraft/Joly files.]}
}
{
}

\Verse{146}
{
    \zA{我闲着还要做老太太屋里的针线，所以宝玉的事竟不曾留 心。}
}
{
    \eA{[English source not present in the checked-in Bencraft/Joly files.]}
}
{
}

\Verse{147}
{
    \zA{太太既怪，从此后我留心就是了。”}
}
{
    \eA{[English source not present in the checked-in Bencraft/Joly files.]}
}
{
}

\Verse{148}
{
    \zA{王夫人信以为实了，忙说：“阿弥陀佛! 你不近宝玉，是我的造化。}
}
{
    \eA{[English source not present in the checked-in Bencraft/Joly files.]}
}
{
}

\Verse{149}
{
    \zA{竟不劳你费心!既是老太太给宝玉的，我明儿回了老太 太再撵你！”}
}
{
    \eA{[English source not present in the checked-in Bencraft/Joly files.]}
}
{
}

\Verse{150}
{
    \zA{因向王善保家的道：“你们进去，好生防他几日，不许他在宝玉屋里 睡觉，等我回过老太太，再处治他。”}
}
{
    \eA{[English source not present in the checked-in Bencraft/Joly files.]}
}
{
}

\Verse{151}
{
    \zA{喝声：“出去!站在这里，我看不上这浪样 儿!谁许你这么花红柳绿的妆扮！”}
}
{
    \eA{[English source not present in the checked-in Bencraft/Joly files.]}
}
{
}

\Verse{152}
{
    \zA{晴雯只得出来。}
}
{
    \eA{[English source not present in the checked-in Bencraft/Joly files.]}
}
{
}

\Verse{153}
{
    \zA{这气非同小可，一出门，便拿 绢子握着脸，一头走，一头哭，直哭到园内去。}
}
{
    \eA{[English source not present in the checked-in Bencraft/Joly files.]}
}
{
}

\Verse{154}
{
    \zA{这里王夫人向凤姐等自怨道：“这几年我越发精神短了，照顾不到，这样妖精 似的东西竟没看见!只怕这样的还有，明日倒得查查。”}
}
{
    \eA{[English source not present in the checked-in Bencraft/Joly files.]}
}
{
}

\Verse{155}
{
    \zA{凤姐见王夫人盛怒之际， 又因王善保家的是邢夫人的耳目，常时调唆的邢夫人生事，纵有千百样言语，此刻 也不敢说，只低头答应着。}
}
{
    \eA{[English source not present in the checked-in Bencraft/Joly files.]}
}
{
}

\Verse{156}
{
    \zA{王善保家的道：“太太且请息怒。}
}
{
    \eA{[English source not present in the checked-in Bencraft/Joly files.]}
}
{
}

\Verse{157}
{
    \zA{这些事小。}
}
{
    \eA{[English source not present in the checked-in Bencraft/Joly files.]}
}
{
}

\Verse{158}
{
    \zA{只交与奴 才。}
}
{
    \eA{[English source not present in the checked-in Bencraft/Joly files.]}
}
{
}

\Verse{159}
{
    \zA{如今要查这个是极容易的。}
}
{
    \eA{[English source not present in the checked-in Bencraft/Joly files.]}
}
{
}

\Verse{160}
{
    \zA{等到晚上园门关了的时节，内外不通风，我们竟给 他们个冷不防，带着人到各处丫头们房里搜寻。}
}
{
    \eA{[English source not present in the checked-in Bencraft/Joly files.]}
}
{
}

\Verse{161}
{
    \zA{想来谁有这个，断不单有这个，自 然还有别的。}
}
{
    \eA{[English source not present in the checked-in Bencraft/Joly files.]}
}
{
}

\Verse{162}
{
    \zA{那时翻出别的来，自然这个也是他的了。”}
}
{
    \eA{[English source not present in the checked-in Bencraft/Joly files.]}
}
{
}

\Verse{163}
{
    \zA{王夫人道：“这话倒是。}
}
{
    \eA{[English source not present in the checked-in Bencraft/Joly files.]}
}
{
}

\Verse{164}
{
    \zA{若不如此，断乎不能明白。”}
}
{
    \eA{[English source not present in the checked-in Bencraft/Joly files.]}
}
{
}

\Verse{165}
{
    \zA{因问凤姐：“如何？”}
}
{
    \eA{[English source not present in the checked-in Bencraft/Joly files.]}
}
{
}

\Verse{166}
{
    \zA{凤姐只得答应说：“太太说是， 就行罢了。”}
}
{
    \eA{[English source not present in the checked-in Bencraft/Joly files.]}
}
{
}

\Verse{167}
{
    \zA{王夫人道：“这主意很是，不然一年也查不出来。”}
}
{
    \eA{[English source not present in the checked-in Bencraft/Joly files.]}
}
{
}

\Verse{168}
{
    \zA{于是大家商议已 定。}
}
{
    \eA{[English source not present in the checked-in Bencraft/Joly files.]}
}
{
}

\Verse{169}
{
    \zA{至晚饭后，待贾母安寝了，宝钗等入园时，王家的便请了凤姐一并进园，喝命 将角门皆上锁，便从上夜的婆子处来抄检起。}
}
{
    \eA{[English source not present in the checked-in Bencraft/Joly files.]}
}
{
}

\Verse{170}
{
    \zA{不过抄检些多馀攒下蜡烛灯油等物。}
}
{
    \eA{[English source not present in the checked-in Bencraft/Joly files.]}
}
{
}

\Verse{171}
{
    \zA{王善保家的道：“这也是赃，不许动的，等明日回过太太再动。”}
}
{
    \eA{[English source not present in the checked-in Bencraft/Joly files.]}
}
{
}

\Verse{172}
{
    \zA{于是先就到怡红 院中，喝命关门。}
}
{
    \eA{[English source not present in the checked-in Bencraft/Joly files.]}
}
{
}

\Verse{173}
{
    \zA{当下宝玉正因晴雯不自在，忽见这一干人来，不知为何直扑了丫 头们的房门去。}
}
{
    \eA{[English source not present in the checked-in Bencraft/Joly files.]}
}
{
}

\Verse{174}
{
    \zA{因迎出凤姐来，问是何故。}
}
{
    \eA{[English source not present in the checked-in Bencraft/Joly files.]}
}
{
}

\Verse{175}
{
    \zA{凤姐道：“丢了一件要紧的东西，因大 家混赖，恐怕有丫头们偷了，所以大家都查一查，去疑儿。”}
}
{
    \eA{[English source not present in the checked-in Bencraft/Joly files.]}
}
{
}

\Verse{176}
{
    \zA{一面说，一面坐下吃 茶。}
}
{
    \eA{[English source not present in the checked-in Bencraft/Joly files.]}
}
{
}

\Verse{177}
{
    \zA{王家的等搜了一回，又细问：“这几个箱子是谁的？”}
}
{
    \eA{[English source not present in the checked-in Bencraft/Joly files.]}
}
{
}

\Verse{178}
{
    \zA{都叫本人来亲自打开。}
}
{
    \eA{[English source not present in the checked-in Bencraft/Joly files.]}
}
{
}

\Verse{179}
{
    \zA{袭人因见晴雯这样，必有异事，又见这番抄检，只得自己先出来打开了箱子并匣子， 任其搜检一番，不过平常通用之物。}
}
{
    \eA{[English source not present in the checked-in Bencraft/Joly files.]}
}
{
}

\Verse{180}
{
    \zA{随放下又搜别人的，挨次都一一搜过。}
}
{
    \eA{[English source not present in the checked-in Bencraft/Joly files.]}
}
{
}

\Verse{181}
{
    \zA{到晴雯 的箱子，因问：“是谁的？}
}
{
    \eA{[English source not present in the checked-in Bencraft/Joly files.]}
}
{
}

\Verse{182}
{
    \zA{怎么不打开叫搜？”}
}
{
    \eA{[English source not present in the checked-in Bencraft/Joly files.]}
}
{
}

\Verse{183}
{
    \zA{袭人方欲替晴雯开时，只见晴雯挽 着头发闯进来，？}
}
{
    \eA{[English source not present in the checked-in Bencraft/Joly files.]}
}
{
}

\Verse{184}
{
    \zA{啷一声将箱子掀开，两手提着底子往地下一倒，将所有之物尽都 倒出来。}
}
{
    \eA{[English source not present in the checked-in Bencraft/Joly files.]}
}
{
}

\Verse{185}
{
    \zA{王善保家的也觉没趣儿，便紫胀了脸，说道：“姑娘你别生气。}
}
{
    \eA{[English source not present in the checked-in Bencraft/Joly files.]}
}
{
}

\Verse{186}
{
    \zA{我们并非 私自就来的，原是奉太太的命来搜察，你们叫翻呢，我们就翻一翻，不叫翻，我们 还许回太太去呢。}
}
{
    \eA{[English source not present in the checked-in Bencraft/Joly files.]}
}
{
}

\Verse{187}
{
    \zA{那用急的这个样子！”}
}
{
    \eA{[English source not present in the checked-in Bencraft/Joly files.]}
}
{
}

\Verse{188}
{
    \zA{晴雯听了这话，越发火上浇油，便指着他 的脸说道：“你说你是太太打发来的，我还是老太太打发来的呢!太太那边的人我
        也都见过，就只没看见你这么个有头有脸大管事的奶奶！”}
}
{
    \eA{[English source not present in the checked-in Bencraft/Joly files.]}
}
{
}

\Verse{189}
{
    \zA{凤姐见晴雯说话锋利尖 酸，心中甚喜，却碍着邢夫人的脸，忙喝住晴雯。}
}
{
    \eA{[English source not present in the checked-in Bencraft/Joly files.]}
}
{
}

\Verse{190}
{
    \zA{那王善保家的又羞又气，刚要还 言，凤姐道：“妈妈，你也不必和他们一般见识，你且细细搜你的，咱们还到各处 走走呢。}
}
{
    \eA{[English source not present in the checked-in Bencraft/Joly files.]}
}
{
}

\Verse{191}
{
    \zA{再迟了走了风，我可担不起。”}
}
{
    \eA{[English source not present in the checked-in Bencraft/Joly files.]}
}
{
}

\Verse{192}
{
    \zA{王善保家的只得咬咬牙，且忍了这口气， 细细的看了一看，也无甚私弊之物。}
}
{
    \eA{[English source not present in the checked-in Bencraft/Joly files.]}
}
{
}

\Verse{193}
{
    \zA{回了凤姐，要别处去，凤姐道：“你可细细的 查，若这一番查不出来，难回话的。”}
}
{
    \eA{[English source not present in the checked-in Bencraft/Joly files.]}
}
{
}

\Verse{194}
{
    \zA{众人都道：“尽都细翻了，没有什么差错东 西。}
}
{
    \eA{[English source not present in the checked-in Bencraft/Joly files.]}
}
{
}

\Verse{195}
{
    \zA{虽有几样男人物件，都是小孩子的东西，想是宝玉的旧物，没甚关系的。”}
}
{
    \eA{[English source not present in the checked-in Bencraft/Joly files.]}
}
{
}

\Verse{196}
{
    \zA{凤 姐听了，笑道：“既如此，咱们就走，再瞧别处去。”}
}
{
    \eA{[English source not present in the checked-in Bencraft/Joly files.]}
}
{
}

\Verse{197}
{
    \zA{说着，一径出来，向王善保家的道：“我有一句话，不知是不是：要抄检只抄 检咱们家的人，薛大姑娘屋里，断乎抄检不得的。”}
}
{
    \eA{[English source not present in the checked-in Bencraft/Joly files.]}
}
{
}

\Verse{198}
{
    \zA{王善保家的笑道：“这个自然， 岂有抄起亲戚家来的。”}
}
{
    \eA{[English source not present in the checked-in Bencraft/Joly files.]}
}
{
}

\Verse{199}
{
    \zA{凤姐点头道：“我也这样说呢。”}
}
{
    \eA{[English source not present in the checked-in Bencraft/Joly files.]}
}
{
}

\Verse{200}
{
    \zA{一头说，一头到了潇湘 馆内。}
}
{
    \eA{[English source not present in the checked-in Bencraft/Joly files.]}
}
{
}

\Verse{201}
{
    \zA{黛玉已睡了，忽报这些人来，不知为甚事。}
}
{
    \eA{[English source not present in the checked-in Bencraft/Joly files.]}
}
{
}

\Verse{202}
{
    \zA{才要起来，只见凤姐已走进来， 忙按住他不叫起来，只说：“睡着罢，我们就走的。”}
}
{
    \eA{[English source not present in the checked-in Bencraft/Joly files.]}
}
{
}

\Verse{203}
{
    \zA{这边且说些闲话。}
}
{
    \eA{[English source not present in the checked-in Bencraft/Joly files.]}
}
{
}

\Verse{204}
{
    \zA{那王善保 家的带了众人到了丫鬟房中，也一一开箱倒笼抄检了一番，因从紫鹃房中搜出两副
        宝玉往常换下来的寄名符儿，一副束带上的帔带，两个荷包并扇套，套内有扇子， 打开看时，皆是宝玉往日手内曾拿过的。}
}
{
    \eA{[English source not present in the checked-in Bencraft/Joly files.]}
}
{
}

\Verse{205}
{
    \zA{王善保家的自为得了意，遂忙请凤姐过来 验视，又说：“这些东西从那里来的？”}
}
{
    \eA{[English source not present in the checked-in Bencraft/Joly files.]}
}
{
}

\Verse{206}
{
    \zA{凤姐笑道：“宝玉和他们从小儿在一处混 了几年，这自然是宝玉的旧东西。}
}
{
    \eA{[English source not present in the checked-in Bencraft/Joly files.]}
}
{
}

\Verse{207}
{
    \zA{况且这符儿合扇子，都是老太太和太太常见的。}
}
{
    \eA{[English source not present in the checked-in Bencraft/Joly files.]}
}
{
}

\Verse{208}
{
    \zA{妈妈不信，咱们只管拿了去。”}
}
{
    \eA{[English source not present in the checked-in Bencraft/Joly files.]}
}
{
}

\Verse{209}
{
    \zA{王家的忙笑道：“二奶奶既知道就是了。”}
}
{
    \eA{[English source not present in the checked-in Bencraft/Joly files.]}
}
{
}

\Verse{210}
{
    \zA{凤姐道： “这也不是什么稀罕事，撂下再往别处去是正经。”}
}
{
    \eA{[English source not present in the checked-in Bencraft/Joly files.]}
}
{
}

\Verse{211}
{
    \zA{紫鹃笑道：“直到如今，我们 两下里的账也算不清，要问这一个，连我也忘了是那年月日有的了。”}
}
{
    \eA{[English source not present in the checked-in Bencraft/Joly files.]}
}
{
}

\Verse{212}
{
    \zA{这里凤姐合王善保家的又到探春院内。}
}
{
    \eA{[English source not present in the checked-in Bencraft/Joly files.]}
}
{
}

\Verse{213}
{
    \zA{谁知早有人报与探春了。}
}
{
    \eA{[English source not present in the checked-in Bencraft/Joly files.]}
}
{
}

\Verse{214}
{
    \zA{探春也就猜着 必有原故，所以引出这等丑态来，遂命众丫鬟秉烛开门而待。}
}
{
    \eA{[English source not present in the checked-in Bencraft/Joly files.]}
}
{
}

\Verse{215}
{
    \zA{一时众人来了，探春 故问：“何事？”}
}
{
    \eA{[English source not present in the checked-in Bencraft/Joly files.]}
}
{
}

\Verse{216}
{
    \zA{凤姐笑道：“因丢了一件东西，连日访察不出人来，恐怕旁人赖 这些女孩子们。}
}
{
    \eA{[English source not present in the checked-in Bencraft/Joly files.]}
}
{
}

\Verse{217}
{
    \zA{所以大家搜一搜，使人去疑儿，倒是洗净他们的好法子。”}
}
{
    \eA{[English source not present in the checked-in Bencraft/Joly files.]}
}
{
}

\Verse{218}
{
    \zA{探春笑 道：“我们的丫头自然都是些贼，我就是头一个窝主。}
}
{
    \eA{[English source not present in the checked-in Bencraft/Joly files.]}
}
{
}

\Verse{219}
{
    \zA{既如此，先来搜我的箱柜， 他们所偷了来的，都交给我藏着呢。”}
}
{
    \eA{[English source not present in the checked-in Bencraft/Joly files.]}
}
{
}

\Verse{220}
{
    \zA{说着，便命丫鬟们把箱一齐打开，将镜奁、 妆盒、衾袱、衣包若大若小之物，一齐打开，请凤姐去抄阅。}
}
{
    \eA{[English source not present in the checked-in Bencraft/Joly files.]}
}
{
}

\Verse{221}
{
    \zA{凤姐陪笑道：“我不 过是奉太太的命来，妹妹别错怪了我。”}
}
{
    \eA{[English source not present in the checked-in Bencraft/Joly files.]}
}
{
}

\Verse{222}
{
    \zA{因命丫鬟们：“快快给姑娘关上。”}
}
{
    \eA{[English source not present in the checked-in Bencraft/Joly files.]}
}
{
}

\Verse{223}
{
    \zA{平儿 丰儿等先忙着替侍书等关的关，收的收。}
}
{
    \eA{[English source not present in the checked-in Bencraft/Joly files.]}
}
{
}

\Verse{224}
{
    \zA{探春道：“我的东西倒许你们搜阅，要想 搜我的丫头这可不能。}
}
{
    \eA{[English source not present in the checked-in Bencraft/Joly files.]}
}
{
}

\Verse{225}
{
    \zA{我原比众人歹毒，凡丫头所有的东西，我都知道，都在我这 里间收着：一针一线，他们也没得收藏。}
}
{
    \eA{[English source not present in the checked-in Bencraft/Joly files.]}
}
{
}

\Verse{226}
{
    \zA{要搜，所以只来搜我。}
}
{
    \eA{[English source not present in the checked-in Bencraft/Joly files.]}
}
{
}

\Verse{227}
{
    \zA{你们不依，只管去 回太太，只说我违背了太太，该怎么处治，我去自领。}
}
{
    \eA{[English source not present in the checked-in Bencraft/Joly files.]}
}
{
}

\Verse{228}
{
    \zA{――你们别忙，自然你们抄 的日子有呢!你们今日早起不是议论甄家，自己盼着好好的抄家，果然今日真抄了!
        咱们也渐渐的来了!可知这样大族人家，若从外头杀来，一时是杀不死的。}
}
{
    \eA{[English source not present in the checked-in Bencraft/Joly files.]}
}
{
}

\Verse{229}
{
    \zA{这可是 古人说的，‘百足之虫，死而不僵’，必须先从家里自杀自灭起来，才能一败涂地 呢！”}
}
{
    \eA{[English source not present in the checked-in Bencraft/Joly files.]}
}
{
}

\Verse{230}
{
    \zA{说着，不觉流下泪来。}
}
{
    \eA{[English source not present in the checked-in Bencraft/Joly files.]}
}
{
}

\Verse{231}
{
    \zA{凤姐只看着众媳妇们。}
}
{
    \eA{[English source not present in the checked-in Bencraft/Joly files.]}
}
{
}

\Verse{232}
{
    \zA{周瑞家的便道：“既是女孩子 的东西全在这里，奶奶且请到别处去罢，也让姑娘好安寝。”}
}
{
    \eA{[English source not present in the checked-in Bencraft/Joly files.]}
}
{
}

\Verse{233}
{
    \zA{凤姐便起身告辞。}
}
{
    \eA{[English source not present in the checked-in Bencraft/Joly files.]}
}
{
}

\Verse{234}
{
    \zA{探 春道：“可细细搜明白了!若明日再来，我就不依了。”}
}
{
    \eA{[English source not present in the checked-in Bencraft/Joly files.]}
}
{
}

\Verse{235}
{
    \zA{凤姐笑道：“既然丫头们 的东西都在这里，就不必搜了。”}
}
{
    \eA{[English source not present in the checked-in Bencraft/Joly files.]}
}
{
}

\Verse{236}
{
    \zA{探春冷笑道：“你果然倒乖!连我的包袱都打开 了，还说没翻，明日敢说我护着丫头们，不许你们翻了。}
}
{
    \eA{[English source not present in the checked-in Bencraft/Joly files.]}
}
{
}

\Verse{237}
{
    \zA{你趁早说明，若还要翻， 不妨再翻一遍。”}
}
{
    \eA{[English source not present in the checked-in Bencraft/Joly files.]}
}
{
}

\Verse{238}
{
    \zA{凤姐知道探春素日与众不同的，只得陪笑道：“已经连你的东西 都搜察明白了。”}
}
{
    \eA{[English source not present in the checked-in Bencraft/Joly files.]}
}
{
}

\Verse{239}
{
    \zA{探春又问众人：“你们也都搜明白了没有？”}
}
{
    \eA{[English source not present in the checked-in Bencraft/Joly files.]}
}
{
}

\Verse{240}
{
    \zA{周瑞家的等都陪笑 说：“都明白了。”}
}
{
    \eA{[English source not present in the checked-in Bencraft/Joly files.]}
}
{
}

\Verse{241}
{
    \zA{那王善保家的本是个心内没成算的人，素日虽闻探春的名，他想众人没眼色、 没胆量罢了，那里一个姑娘就这样利害起来？}
}
{
    \eA{[English source not present in the checked-in Bencraft/Joly files.]}
}
{
}

\Verse{242}
{
    \zA{况且又是庶出，他敢怎么着？}
}
{
    \eA{[English source not present in the checked-in Bencraft/Joly files.]}
}
{
}

\Verse{243}
{
    \zA{自己又仗 着是邢夫人的陪房，连王夫人尚另眼相待，何况别人？}
}
{
    \eA{[English source not present in the checked-in Bencraft/Joly files.]}
}
{
}

\Verse{244}
{
    \zA{只当是探春认真单恼凤姐， 与他们无干。}
}
{
    \eA{[English source not present in the checked-in Bencraft/Joly files.]}
}
{
}

\Verse{245}
{
    \zA{他便要趁势作脸，因越众向前，拉起探春的衣襟，故意一掀，嘻嘻的 笑道：“连姑娘身上我都翻了，果然没有什么。”}
}
{
    \eA{[English source not present in the checked-in Bencraft/Joly files.]}
}
{
}

\Verse{246}
{
    \zA{凤姐见他这样，忙说：“妈妈走 罢，别疯疯癫癫的――”一语未了，只听“拍”的一声，王家的脸上早着了探春一 巴掌。}
}
{
    \eA{[English source not present in the checked-in Bencraft/Joly files.]}
}
{
}

\Verse{247}
{
    \zA{探春登时大怒，指着王家的问道：“你是什么东西，敢来拉扯我的衣裳!我 不过看着太太的面上，你又有几岁年纪，叫你一声‘妈妈’，你就狗仗人势，天天
        作耗，在我们跟前逞脸。}
}
{
    \eA{[English source not present in the checked-in Bencraft/Joly files.]}
}
{
}

\Verse{248}
{
    \zA{如今越发了不得了，你索性望我动手动脚的了!你打量我 是和你们姑娘那么好性儿，由着你们欺负？}
}
{
    \eA{[English source not present in the checked-in Bencraft/Joly files.]}
}
{
}

\Verse{249}
{
    \zA{你就错了主意了!你来搜检东西我不恼， 你不该拿我取笑儿！”}
}
{
    \eA{[English source not present in the checked-in Bencraft/Joly files.]}
}
{
}

\Verse{250}
{
    \zA{说着，便亲自要解钮子，拉着凤姐儿细细的翻，“省得叫你 们奴才来翻我！”}
}
{
    \eA{[English source not present in the checked-in Bencraft/Joly files.]}
}
{
}

\Verse{251}
{
    \zA{凤姐平儿等都忙与探春理裙整袂，口内喝着王善保家的说：“妈妈吃两口酒， 就疯疯癫癫起来，前儿把太太也冲撞了。}
}
{
    \eA{[English source not present in the checked-in Bencraft/Joly files.]}
}
{
}

\Verse{252}
{
    \zA{快出去，别再讨脸了！”}
}
{
    \eA{[English source not present in the checked-in Bencraft/Joly files.]}
}
{
}

\Verse{253}
{
    \zA{又忙劝探春：“好 姑娘，别生气。}
}
{
    \eA{[English source not present in the checked-in Bencraft/Joly files.]}
}
{
}

\Verse{254}
{
    \zA{他算什么，姑娘气着倒值多了。”}
}
{
    \eA{[English source not present in the checked-in Bencraft/Joly files.]}
}
{
}

\Verse{255}
{
    \zA{探春冷笑道：“我但凡有气，早 一头碰死了。}
}
{
    \eA{[English source not present in the checked-in Bencraft/Joly files.]}
}
{
}

\Verse{256}
{
    \zA{不然，怎么许奴才来我身上搜贼赃呢!明儿一早，先回过老太太、太 太，再过去给大娘赔礼。}
}
{
    \eA{[English source not present in the checked-in Bencraft/Joly files.]}
}
{
}

\Verse{257}
{
    \zA{该怎么着，我去领！”}
}
{
    \eA{[English source not present in the checked-in Bencraft/Joly files.]}
}
{
}

\Verse{258}
{
    \zA{那王善保家的讨了个没脸，赶忙躲 出窗外，只说：“罢了，罢了!这也是头一遭挨打!我明儿回了太太，仍回老娘家去 罢，这个老命还要他做什么。”}
}
{
    \eA{[English source not present in the checked-in Bencraft/Joly files.]}
}
{
}

\Verse{259}
{
    \zA{探春喝命丫鬟：“你们听着他说话，还等我和他拌 嘴去不成？”}
}
{
    \eA{[English source not present in the checked-in Bencraft/Joly files.]}
}
{
}

\Verse{260}
{
    \zA{侍书听说，便出去说道：“妈妈，你知点道理儿，省一句儿罢。}
}
{
    \eA{[English source not present in the checked-in Bencraft/Joly files.]}
}
{
}

\Verse{261}
{
    \zA{你果 然回老娘家去，倒是我们的造化了，只怕你舍不得去。}
}
{
    \eA{[English source not present in the checked-in Bencraft/Joly files.]}
}
{
}

\Verse{262}
{
    \zA{你去了，叫谁讨主子的好儿， 调唆着察考姑娘、折磨我们呢？”}
}
{
    \eA{[English source not present in the checked-in Bencraft/Joly files.]}
}
{
}

\Verse{263}
{
    \zA{凤姐笑道：“好丫头，真是有其主必有其仆。”}
}
{
    \eA{[English source not present in the checked-in Bencraft/Joly files.]}
}
{
}

\Verse{264}
{
    \zA{探春冷笑道：“我们做贼的人，嘴里都有三言两语的，就只不会背地里调唆主子！”}
}
{
    \eA{[English source not present in the checked-in Bencraft/Joly files.]}
}
{
}

\Verse{265}
{
    \zA{平儿忙也陪笑解劝，一面又拉了侍书进来。}
}
{
    \eA{[English source not present in the checked-in Bencraft/Joly files.]}
}
{
}

\Verse{266}
{
    \zA{周瑞家的等人劝了一番，凤姐直待伏侍 探春睡下，方带着人往对过暖香坞来。}
}
{
    \eA{[English source not present in the checked-in Bencraft/Joly files.]}
}
{
}

\Verse{267}
{
    \zA{彼时李纨犹病在床上，他与惜春是紧邻，又和探春相近，故顺路先到这两处。}
}
{
    \eA{[English source not present in the checked-in Bencraft/Joly files.]}
}
{
}

\Verse{268}
{
    \zA{因李纨才吃了药睡着，不好惊动，只到丫鬟们房中，一一的搜了一遍，也没有什么 东西，遂到惜春房中来。}
}
{
    \eA{[English source not present in the checked-in Bencraft/Joly files.]}
}
{
}

\Verse{269}
{
    \zA{因惜春年少，尚未识事，吓的不知当有什么事故，凤姐少 不得安慰他。}
}
{
    \eA{[English source not present in the checked-in Bencraft/Joly files.]}
}
{
}

\Verse{270}
{
    \zA{谁知竟在入画箱中寻出一大包银锞子来，约共三四十个，为察奸情， 反得贼赃。}
}
{
    \eA{[English source not present in the checked-in Bencraft/Joly files.]}
}
{
}

\Verse{271}
{
    \zA{又有一副玉带版子，并一包男人的靴袜等物。}
}
{
    \eA{[English source not present in the checked-in Bencraft/Joly files.]}
}
{
}

\Verse{272}
{
    \zA{凤姐也黄了脸，因问：“是 那里来的？”}
}
{
    \eA{[English source not present in the checked-in Bencraft/Joly files.]}
}
{
}

\Verse{273}
{
    \zA{入画只得跪下哭诉真情，说：“这是珍大爷赏我哥哥的。}
}
{
    \eA{[English source not present in the checked-in Bencraft/Joly files.]}
}
{
}

\Verse{274}
{
    \zA{因我们老子 娘都在南方，如今只跟着叔叔过日子；我叔叔婶子只要喝酒赌钱，我哥哥怕交给他 们又花了，所以每常得了，悄悄的烦老妈妈带进来，叫我收着的。”}
}
{
    \eA{[English source not present in the checked-in Bencraft/Joly files.]}
}
{
}

\Verse{275}
{
    \zA{惜春胆小，见 了这个，也害怕说：“我竟不知道，这还了得。}
}
{
    \eA{[English source not present in the checked-in Bencraft/Joly files.]}
}
{
}

\Verse{276}
{
    \zA{二嫂子要打他，好歹带出他去打罢， 我听不惯的。”}
}
{
    \eA{[English source not present in the checked-in Bencraft/Joly files.]}
}
{
}

\Verse{277}
{
    \zA{凤姐笑道：“若果真呢，也倒可恕，只是不该私自传送进来。}
}
{
    \eA{[English source not present in the checked-in Bencraft/Joly files.]}
}
{
}

\Verse{278}
{
    \zA{这个 可以传递，怕什么不可传递？}
}
{
    \eA{[English source not present in the checked-in Bencraft/Joly files.]}
}
{
}

\Verse{279}
{
    \zA{这倒是传递人的不是了。}
}
{
    \eA{[English source not present in the checked-in Bencraft/Joly files.]}
}
{
}

\Verse{280}
{
    \zA{若这话不真，倘是偷来的， 你可就别想活了。”}
}
{
    \eA{[English source not present in the checked-in Bencraft/Joly files.]}
}
{
}

\Verse{281}
{
    \zA{入画跪哭道：“我不敢撒谎，奶奶只管明日问我们奶奶和大爷 去，若说不是赏的，就拿我和我哥哥一同打死无怨。”}
}
{
    \eA{[English source not present in the checked-in Bencraft/Joly files.]}
}
{
}

\Verse{282}
{
    \zA{凤姐道：“这个自然要问的。}
}
{
    \eA{[English source not present in the checked-in Bencraft/Joly files.]}
}
{
}

\Verse{283}
{
    \zA{只是真赏的，也有不是，谁许你私自传送东西呢？}
}
{
    \eA{[English source not present in the checked-in Bencraft/Joly files.]}
}
{
}

\Verse{284}
{
    \zA{你且说是谁接的，我就饶你。}
}
{
    \eA{[English source not present in the checked-in Bencraft/Joly files.]}
}
{
}

\Verse{285}
{
    \zA{下 次万万不可。”}
}
{
    \eA{[English source not present in the checked-in Bencraft/Joly files.]}
}
{
}

\Verse{286}
{
    \zA{惜春道：“嫂子别饶他，这里人多，要不管了他，那些大的听见了 又不知怎么样呢。}
}
{
    \eA{[English source not present in the checked-in Bencraft/Joly files.]}
}
{
}

\Verse{287}
{
    \zA{嫂子要依他，我也不依。”}
}
{
    \eA{[English source not present in the checked-in Bencraft/Joly files.]}
}
{
}

\Verse{288}
{
    \zA{凤姐道：“素日我看他还使得。}
}
{
    \eA{[English source not present in the checked-in Bencraft/Joly files.]}
}
{
}

\Verse{289}
{
    \zA{谁没 一个错，只这一次。}
}
{
    \eA{[English source not present in the checked-in Bencraft/Joly files.]}
}
{
}

\Verse{290}
{
    \zA{二次再犯，两罪俱罚。}
}
{
    \eA{[English source not present in the checked-in Bencraft/Joly files.]}
}
{
}

\Verse{291}
{
    \zA{但不知传递是谁？”}
}
{
    \eA{[English source not present in the checked-in Bencraft/Joly files.]}
}
{
}

\Verse{292}
{
    \zA{惜春道：“若说传 递，再无别人，必是后门上的老张。}
}
{
    \eA{[English source not present in the checked-in Bencraft/Joly files.]}
}
{
}

\Verse{293}
{
    \zA{他常和这些丫头们鬼鬼祟祟的，这些丫头们也 都肯照顾他。”}
}
{
    \eA{[English source not present in the checked-in Bencraft/Joly files.]}
}
{
}

\Verse{294}
{
    \zA{凤姐听说，便命人记下，将东西且交给周瑞家的暂且拿着，等明日 对明再议。}
}
{
    \eA{[English source not present in the checked-in Bencraft/Joly files.]}
}
{
}

\Verse{295}
{
    \zA{谁知那老张妈原和王善保家有亲，近因王善保家的在邢夫人跟前作了心 腹人，便把亲戚和伴儿们都看不到眼里了。}
}
{
    \eA{[English source not present in the checked-in Bencraft/Joly files.]}
}
{
}

\Verse{296}
{
    \zA{后来张家的气不平，斗了两次口，彼此 都不说话了。}
}
{
    \eA{[English source not present in the checked-in Bencraft/Joly files.]}
}
{
}

\Verse{297}
{
    \zA{如今王家的听见是他传递，碰在他心坎儿上，更兼刚才挨了探春的打， 受了侍书的气，没处发泄，听见张家的这事，因撺掇凤姐道：“这传东西的事关系 更大。}
}
{
    \eA{[English source not present in the checked-in Bencraft/Joly files.]}
}
{
}

\Verse{298}
{
    \zA{想来那些东西，自然也是传递进来的。}
}
{
    \eA{[English source not present in the checked-in Bencraft/Joly files.]}
}
{
}

\Verse{299}
{
    \zA{奶奶倒不可不问。”}
}
{
    \eA{[English source not present in the checked-in Bencraft/Joly files.]}
}
{
}

\Verse{300}
{
    \zA{凤姐儿道：“我 知道，不用你说。”}
}
{
    \eA{[English source not present in the checked-in Bencraft/Joly files.]}
}
{
}

\Verse{301}
{
    \zA{于是别了惜春，方往迎春房内去。}
}
{
    \eA{[English source not present in the checked-in Bencraft/Joly files.]}
}
{
}

\Verse{302}
{
    \zA{迎春已经睡着了，丫鬟们也才要睡，众人扣 门，半日才开。}
}
{
    \eA{[English source not present in the checked-in Bencraft/Joly files.]}
}
{
}

\Verse{303}
{
    \zA{凤姐吩咐：“不必惊动姑娘。”}
}
{
    \eA{[English source not present in the checked-in Bencraft/Joly files.]}
}
{
}

\Verse{304}
{
    \zA{遂往丫鬟们房里来。}
}
{
    \eA{[English source not present in the checked-in Bencraft/Joly files.]}
}
{
}

\Verse{305}
{
    \zA{因司棋是王善 保家的外孙女儿，凤姐要看王家的可藏私不藏，遂留神看他搜检。}
}
{
    \eA{[English source not present in the checked-in Bencraft/Joly files.]}
}
{
}

\Verse{306}
{
    \zA{先从别人箱子搜 起，皆无别物。}
}
{
    \eA{[English source not present in the checked-in Bencraft/Joly files.]}
}
{
}

\Verse{307}
{
    \zA{及到了司棋箱中，随意掏了一回，王善保家的说：“也没有什么东 西。”}
}
{
    \eA{[English source not present in the checked-in Bencraft/Joly files.]}
}
{
}

\Verse{308}
{
    \zA{才要关箱时，周瑞家的道：“这是什么话？}
}
{
    \eA{[English source not present in the checked-in Bencraft/Joly files.]}
}
{
}

\Verse{309}
{
    \zA{有没有，总要一样看看才公道。”}
}
{
    \eA{[English source not present in the checked-in Bencraft/Joly files.]}
}
{
}

\Verse{310}
{
    \zA{说着，便伸手掣出一双男子的绵袜并一双缎鞋，又有一个小包袱。}
}
{
    \eA{[English source not present in the checked-in Bencraft/Joly files.]}
}
{
}

\Verse{311}
{
    \zA{打开看时，里面 是一个同心如意，并一个字帖儿。}
}
{
    \eA{[English source not present in the checked-in Bencraft/Joly files.]}
}
{
}

\Verse{312}
{
    \zA{一总递给凤姐。}
}
{
    \eA{[English source not present in the checked-in Bencraft/Joly files.]}
}
{
}

\Verse{313}
{
    \zA{凤姐因理家久了，每每看帖看帐， 也颇识得几个字了。}
}
{
    \eA{[English source not present in the checked-in Bencraft/Joly files.]}
}
{
}

\Verse{314}
{
    \zA{那帖是大红双喜笺，便看上面写道： 上月你来家后，父母已觉察了。}
}
{
    \eA{[English source not present in the checked-in Bencraft/Joly files.]}
}
{
}

\Verse{315}
{
    \zA{但姑娘未出阁，尚不能完你我心愿。}
}
{
    \eA{[English source not present in the checked-in Bencraft/Joly files.]}
}
{
}

\Verse{316}
{
    \zA{若园内可 以相见，你可托张妈给一信。}
}
{
    \eA{[English source not present in the checked-in Bencraft/Joly files.]}
}
{
}

\Verse{317}
{
    \zA{若得在园内一见，倒比来家好说话。}
}
{
    \eA{[English source not present in the checked-in Bencraft/Joly files.]}
}
{
}

\Verse{318}
{
    \zA{千万千万!再所 赐香珠二串，今已查收。}
}
{
    \eA{[English source not present in the checked-in Bencraft/Joly files.]}
}
{
}

\Verse{319}
{
    \zA{外特寄香袋一个，略表我心。}
}
{
    \eA{[English source not present in the checked-in Bencraft/Joly files.]}
}
{
}

\Verse{320}
{
    \zA{千万收好。}
}
{
    \eA{[English source not present in the checked-in Bencraft/Joly files.]}
}
{
}

\Verse{321}
{
    \zA{表弟潘又安具。}
}
{
    \eA{[English source not present in the checked-in Bencraft/Joly files.]}
}
{
}

\Verse{322}
{
    \zA{凤姐看了，不由的笑将起来。}
}
{
    \eA{[English source not present in the checked-in Bencraft/Joly files.]}
}
{
}

\Verse{323}
{
    \zA{那王善保家的素日并不知道他姑表兄妹有这一节风流 故事，见了这鞋袜，心内已有些毛病，又见有一红帖，凤姐看着笑，他便说道：“必
        是他们写的帐不成字，所以奶奶见笑。”}
}
{
    \eA{[English source not present in the checked-in Bencraft/Joly files.]}
}
{
}

\Verse{324}
{
    \zA{凤姐笑道：“正是这个帐竟算不过来!你 是司棋的老娘，你表弟也该姓王，怎么又姓潘呢？”}
}
{
    \eA{[English source not present in the checked-in Bencraft/Joly files.]}
}
{
}

\Verse{325}
{
    \zA{王善保家的见问的奇怪，只得 勉强告道：“司棋的姑妈给了潘家，所以他姑表弟兄姓潘。}
}
{
    \eA{[English source not present in the checked-in Bencraft/Joly files.]}
}
{
}

\Verse{326}
{
    \zA{上次逃走了的潘又安， 就是他。”}
}
{
    \eA{[English source not present in the checked-in Bencraft/Joly files.]}
}
{
}

\Verse{327}
{
    \zA{凤姐笑道：“这就是了。”}
}
{
    \eA{[English source not present in the checked-in Bencraft/Joly files.]}
}
{
}

\Verse{328}
{
    \zA{因说：“我念给你听听。”}
}
{
    \eA{[English source not present in the checked-in Bencraft/Joly files.]}
}
{
}

\Verse{329}
{
    \zA{说着，从头念了 一遍，大家都吓一跳。}
}
{
    \eA{[English source not present in the checked-in Bencraft/Joly files.]}
}
{
}

\Verse{330}
{
    \zA{这王家的一心只要拿人的错儿，不想反拿住了他外孙女儿， 又气又臊。}
}
{
    \eA{[English source not present in the checked-in Bencraft/Joly files.]}
}
{
}

\Verse{331}
{
    \zA{周瑞家的四人听见凤姐儿念了，都吐舌头，摇头儿。}
}
{
    \eA{[English source not present in the checked-in Bencraft/Joly files.]}
}
{
}

\Verse{332}
{
    \zA{周瑞家的道：“王 大妈听见了!这是明明白白，再没得话说了。}
}
{
    \eA{[English source not present in the checked-in Bencraft/Joly files.]}
}
{
}

\Verse{333}
{
    \zA{这如今怎么样呢？”}
}
{
    \eA{[English source not present in the checked-in Bencraft/Joly files.]}
}
{
}

\Verse{334}
{
    \zA{王家的只恨无地 缝儿可钻。}
}
{
    \eA{[English source not present in the checked-in Bencraft/Joly files.]}
}
{
}

\Verse{335}
{
    \zA{凤姐只瞅着他，抿着嘴儿嘻嘻的笑，向周瑞家的道：“这倒也好。}
}
{
    \eA{[English source not present in the checked-in Bencraft/Joly files.]}
}
{
}

\Verse{336}
{
    \zA{不用 他老娘操一点心儿，鸦雀不闻，就给他们弄了个好女婿来了。”}
}
{
    \eA{[English source not present in the checked-in Bencraft/Joly files.]}
}
{
}

\Verse{337}
{
    \zA{周瑞家的也笑着凑 趣儿。}
}
{
    \eA{[English source not present in the checked-in Bencraft/Joly files.]}
}
{
}

\Verse{338}
{
    \zA{王家的无处煞气，只好打着自己的脸骂道：“老不死的娼妇，怎么造下孽了？}
}
{
    \eA{[English source not present in the checked-in Bencraft/Joly files.]}
}
{
}

\Verse{339}
{
    \zA{说嘴打嘴，现世现报！”}
}
{
    \eA{[English source not present in the checked-in Bencraft/Joly files.]}
}
{
}

\Verse{340}
{
    \zA{众人见他如此，要笑又不敢笑，也有趁愿的，也有心中感 动报应不爽的。}
}
{
    \eA{[English source not present in the checked-in Bencraft/Joly files.]}
}
{
}

\Verse{341}
{
    \zA{凤姐见司棋低头不语，也并无畏惧惭愧之意，倒觉可异。}
}
{
    \eA{[English source not present in the checked-in Bencraft/Joly files.]}
}
{
}

\Verse{342}
{
    \zA{料此时夜深，且不必 盘问，只怕他夜间自寻短志，遂唤两个婆子监守，且带了人，拿了赃证，回来歇息， 等待明日料理。}
}
{
    \eA{[English source not present in the checked-in Bencraft/Joly files.]}
}
{
}

\Verse{343}
{
    \zA{谁知夜里下面淋血不止，次日便觉身体十分软弱起来，遂掌不住， 请医诊视；开方立案，说要保重而去。}
}
{
    \eA{[English source not present in the checked-in Bencraft/Joly files.]}
}
{
}

\Verse{344}
{
    \zA{老嬷嬷们拿了方子，回过王夫人，不免又添 一番愁闷，遂将司棋之事暂且搁起。}
}
{
    \eA{[English source not present in the checked-in Bencraft/Joly files.]}
}
{
}

\Verse{345}
{
    \zA{可巧这日尤氏来看凤姐，坐了一回，又看李纨等。}
}
{
    \eA{[English source not present in the checked-in Bencraft/Joly files.]}
}
{
}

\Verse{346}
{
    \zA{忽见惜春遣人来请，尤氏到 他房中，惜春便将昨夜之事细细告诉了，又命人将入画的东西一概要来与尤氏过 目。}
}
{
    \eA{[English source not present in the checked-in Bencraft/Joly files.]}
}
{
}

\Verse{347}
{
    \zA{尤氏道：“实是你哥哥赏他哥哥的。}
}
{
    \eA{[English source not present in the checked-in Bencraft/Joly files.]}
}
{
}

\Verse{348}
{
    \zA{只不该私自传送，如今官盐反成了私盐了。”}
}
{
    \eA{[English source not present in the checked-in Bencraft/Joly files.]}
}
{
}

\Verse{349}
{
    \zA{因骂入画：“糊涂东西！”}
}
{
    \eA{[English source not present in the checked-in Bencraft/Joly files.]}
}
{
}

\Verse{350}
{
    \zA{惜春道：“你们管教不严，反骂丫头。}
}
{
    \eA{[English source not present in the checked-in Bencraft/Joly files.]}
}
{
}

\Verse{351}
{
    \zA{这些姊妹，独我 的丫头没脸，我如何去见人!昨儿叫凤姐姐带了他去，又不肯。}
}
{
    \eA{[English source not present in the checked-in Bencraft/Joly files.]}
}
{
}

\Verse{352}
{
    \zA{今日嫂子来的恰好， 快带了他去，或打或杀或卖，我一概不管。”}
}
{
    \eA{[English source not present in the checked-in Bencraft/Joly files.]}
}
{
}

\Verse{353}
{
    \zA{入画听说，跪地哀求，百般苦告。}
}
{
    \eA{[English source not present in the checked-in Bencraft/Joly files.]}
}
{
}

\Verse{354}
{
    \zA{尤 氏和奶妈等人也都十分解说：“他不过一时糊涂，下次再不敢的。}
}
{
    \eA{[English source not present in the checked-in Bencraft/Joly files.]}
}
{
}

\Verse{355}
{
    \zA{看他从小儿伏侍 一场。”}
}
{
    \eA{[English source not present in the checked-in Bencraft/Joly files.]}
}
{
}

\Verse{356}
{
    \zA{谁知惜春年幼，天性孤僻，任人怎说，只是咬定牙，断乎不肯留着。}
}
{
    \eA{[English source not present in the checked-in Bencraft/Joly files.]}
}
{
}

\Verse{357}
{
    \zA{更又 说道：“不但不要入画，如今我也大了，连我也不便往你们那边去了。}
}
{
    \eA{[English source not present in the checked-in Bencraft/Joly files.]}
}
{
}

\Verse{358}
{
    \zA{况且近日闻 得多少议论，我若再去，连我也编派。”}
}
{
    \eA{[English source not present in the checked-in Bencraft/Joly files.]}
}
{
}

\Verse{359}
{
    \zA{尤氏道：“谁敢议论什么？}
}
{
    \eA{[English source not present in the checked-in Bencraft/Joly files.]}
}
{
}

\Verse{360}
{
    \zA{又有什么可议 论的？}
}
{
    \eA{[English source not present in the checked-in Bencraft/Joly files.]}
}
{
}

\Verse{361}
{
    \zA{姑娘是谁？}
}
{
    \eA{[English source not present in the checked-in Bencraft/Joly files.]}
}
{
}

\Verse{362}
{
    \zA{我们是谁？}
}
{
    \eA{[English source not present in the checked-in Bencraft/Joly files.]}
}
{
}

\Verse{363}
{
    \zA{姑娘既听见人议论我们，就该问着他才是。”}
}
{
    \eA{[English source not present in the checked-in Bencraft/Joly files.]}
}
{
}

\Verse{364}
{
    \zA{惜春冷笑 道：“你这话问着我倒好!我一个姑娘家，只好躲是非的，我反寻是非，成个什么 人了。}
}
{
    \eA{[English source not present in the checked-in Bencraft/Joly files.]}
}
{
}

\Verse{365}
{
    \zA{况且古人说的，‘善恶生死，父子不能有所勖助’，何况你我二人之间。}
}
{
    \eA{[English source not present in the checked-in Bencraft/Joly files.]}
}
{
}

\Verse{366}
{
    \zA{我 只能保住自己就够了，以后你们有事好歹别累我。”}
}
{
    \eA{[English source not present in the checked-in Bencraft/Joly files.]}
}
{
}

\Verse{367}
{
    \zA{尤氏听了，又气又好笑，因向 地下众人道：“怪道人人都说四姑娘年轻糊涂，我只不信。}
}
{
    \eA{[English source not present in the checked-in Bencraft/Joly files.]}
}
{
}

\Verse{368}
{
    \zA{你们听这些话，无原无 故，又没轻重，真真的叫人寒心。”}
}
{
    \eA{[English source not present in the checked-in Bencraft/Joly files.]}
}
{
}

\Verse{369}
{
    \zA{众人都劝说道：“姑娘年轻，奶奶自然该吃些 亏的。”}
}
{
    \eA{[English source not present in the checked-in Bencraft/Joly files.]}
}
{
}

\Verse{370}
{
    \zA{惜春冷笑道：“我虽年轻，这话却不年轻。}
}
{
    \eA{[English source not present in the checked-in Bencraft/Joly files.]}
}
{
}

\Verse{371}
{
    \zA{你们不看书，不识字，所以都 是呆子，倒说我糊涂。”}
}
{
    \eA{[English source not present in the checked-in Bencraft/Joly files.]}
}
{
}

\Verse{372}
{
    \zA{尤氏道：“你是状元，第一个才子!我们糊涂人，不如你 明白。”}
}
{
    \eA{[English source not present in the checked-in Bencraft/Joly files.]}
}
{
}

\Verse{373}
{
    \zA{惜春道：“据你这话就不明白。}
}
{
    \eA{[English source not present in the checked-in Bencraft/Joly files.]}
}
{
}

\Verse{374}
{
    \zA{状元难道没有糊涂的？}
}
{
    \eA{[English source not present in the checked-in Bencraft/Joly files.]}
}
{
}

\Verse{375}
{
    \zA{可知你们这些人都 是世俗之见，那里眼里识的出真假、心里分的出好歹来？}
}
{
    \eA{[English source not present in the checked-in Bencraft/Joly files.]}
}
{
}

\Verse{376}
{
    \zA{你们要看真人，总在最初 一步的心上看起，才能明白呢。”}
}
{
    \eA{[English source not present in the checked-in Bencraft/Joly files.]}
}
{
}

\Verse{377}
{
    \zA{尤氏笑道：“好，好，才是才子，这会子又做大 和尚，讲起参悟来了。”}
}
{
    \eA{[English source not present in the checked-in Bencraft/Joly files.]}
}
{
}

\Verse{378}
{
    \zA{惜春道：“我也不是什么参悟。}
}
{
    \eA{[English source not present in the checked-in Bencraft/Joly files.]}
}
{
}

\Verse{379}
{
    \zA{我看如今人一概也都是入 画一般，没有什么大说头儿。”}
}
{
    \eA{[English source not present in the checked-in Bencraft/Joly files.]}
}
{
}

\Verse{380}
{
    \zA{尤氏道：“可知你真是个心冷嘴冷的人。”}
}
{
    \eA{[English source not present in the checked-in Bencraft/Joly files.]}
}
{
}

\Verse{381}
{
    \zA{惜春道： “怎么我不冷!我清清白白的一个人，为什么叫你们带累坏了？”}
}
{
    \eA{[English source not present in the checked-in Bencraft/Joly files.]}
}
{
}

\Verse{382}
{
    \zA{尤氏心内原有病，怕说这些话，听说有人议论，已是心中羞恼，只是今日惜春 分中不好发作，忍耐了大半天。}
}
{
    \eA{[English source not present in the checked-in Bencraft/Joly files.]}
}
{
}

\Verse{383}
{
    \zA{今见惜春又说这话，因按捺不住，便问道：“怎么 就带累了你？}
}
{
    \eA{[English source not present in the checked-in Bencraft/Joly files.]}
}
{
}

\Verse{384}
{
    \zA{你的丫头的不是，无故说我；我倒忍了这半日，你倒越发得了意，只 管说这些话。}
}
{
    \eA{[English source not present in the checked-in Bencraft/Joly files.]}
}
{
}

\Verse{385}
{
    \zA{你是千金小姐，我们以后就不亲近你，仔细带累了小姐的美名儿!即 刻就叫人将入画带了过去。”}
}
{
    \eA{[English source not present in the checked-in Bencraft/Joly files.]}
}
{
}

\Verse{386}
{
    \zA{说着，便赌气起身去了。}
}
{
    \eA{[English source not present in the checked-in Bencraft/Joly files.]}
}
{
}

\Verse{387}
{
    \zA{惜春道：“你这一去了，若 果然不来，倒也省了口舌是非，大家倒还干净。”}
}
{
    \eA{[English source not present in the checked-in Bencraft/Joly files.]}
}
{
}

\Verse{388}
{
    \zA{尤氏听了，越发生气，但终久他 是姑娘，任凭怎么样也不好和他认真的拌起嘴来，只得索性忍了这口气。}
}
{
    \eA{[English source not present in the checked-in Bencraft/Joly files.]}
}
{
}

\Verse{389}
{
    \zA{便也不答 言，一径往前边去了。}
}
{
    \eA{[English source not present in the checked-in Bencraft/Joly files.]}
}
{
}

\Verse{390}
{
    \zA{未知后事如何，且听下回分解。}
}
{
    \eA{[English source not present in the checked-in Bencraft/Joly files.]}
}
{
}

\Verse{391}
{
    \zA{(本章完)}
}
{
    \eA{[English source not present in the checked-in Bencraft/Joly files.]}
}
{
}
